%==========================================================================================
\section{\texttt{bliss\_hdf5\_reader}: reading BLISS scans}\label{sec:bliss}
%==========================================================================================

\subsection{Overview}

BLISS, the ESRF beamline control system, writes one HDF5 \emph{master file} per dataset,
organised along NeXus conventions \citep{Konnecke2015}. Every scan is an entry named
\fn{"<number>.<subscan>"} (\fn{"12.1"}) containing:
\begin{itemize}[nosep]
\item \fn{title}: the command as typed (\fn{"amesh xech 0 1 50 yech 0 1 40 0.1"}),
  \fn{start_time}, \fn{end_time}, \fn{end_reason} (\fn{SUCCESS}, \fn{USER_ABORT}, \dots);
\item \fn{measurement/}: one dataset per recorded quantity, one entry per scan point along
  the first axis (scalar counters and motor positions: shape $(n)$; MCA spectra:
  $(n, n_\mathrm{ch})$; detector images: $(n, h, w)$). These are links to the
  \fn{instrument/} group;
\item \fn{instrument/positioners_start}: every motor position at the start of the scan;
  \fn{instrument/<name>} groups with an \fn{NX_class} attribute (\fn{NXpositioner} for motors);
\item \fn{scan_parameters/}: \fn{npoints}, \fn{count_time} and other parameters, when recorded;
\item \fn{writer/status}, \fn{sample/name}.
\end{itemize}
Detector images are written by Lima in separate files of a \fn{scanNNNN/} folder next to the
master file; the master file exposes them as an HDF5 \emph{virtual dataset}, which maps
ranges of the virtual array onto datasets of the source files.

The module turns this into a \fn{BlissScan} object that knows the \emph{geometry} of the scan
(its shape in acquisition order, the motor driving each dimension, back-and-forth rows,
aborted points) and reads the data on demand. The geometry comes from parsing the command in
the title when it is recognised (Table~\ref{tab:bliss-titles}), and otherwise from the
recorded motor positions, so that scans run through user macros with custom titles are
still read correctly.

\begin{table}[h]
\centering\small
\caption{Scan commands recognised in titles, and the resulting shape (slowest axis first).}
\label{tab:bliss-titles}
\begin{tabularx}{\textwidth}{@{}l X l@{}}
\toprule
Command & Title grammar & Shape \\
\midrule
\fn{ascan dscan a2scan d3scan lineup} & $(\text{motor}\ \text{start}\ \text{stop})\times N$,
intervals, count time & $(\text{intervals}+1,)$ \\
\fn{amesh dmesh} & fast motor start stop intervals, slow motor start stop intervals, time
(the \emph{first} motor is the fast one) & $(n_\mathrm{slow}, n_\mathrm{fast})$ \\
\fn{fscan fscan2d fscan3d} & $(\text{motor}\ \text{start}\ \text{step}\ \text{npoints})\times
d$, acquisition time [period] (the \emph{last} motor is the fast one) & $(n_0,\dots,n_{d-1})$ \\
\texttt{f2scan f3scan \dots} & $(\text{motor}\ \text{start}\ \text{step})\times N$, npoints,
time & $(n,)$ \\
\fn{loopscan} & npoints, time & $(n,)$ \\
\fn{timescan ftimescan} & time [npoints] (0 or absent: until stopped) & $(n,)$ \\
\fn{ct sct} & [time] & $()$ \\
\bottomrule
\end{tabularx}
\end{table}

\begin{hyp}
\begin{itemize}[nosep,leftmargin=1em]
\item Data are stored in acquisition order, one entry per point along the first axis.
  Meshes are row-major with the slow motor first.
\item The number of recorded points is the most common first-axis length of the measurement
  datasets (fast scans add ``\fn{*_raw}'' channels with $2n+1$ entries, which are ignored).
\item A mesh scanned back and forth (snake) is recognised when the second row of the fast
  motor matches the reversed first row better than the first row itself.
\item An aborted 1D scan keeps its recorded length; an aborted mesh keeps its planned shape
  and missing points are NaN in gridded data.
\item Frames are read through the master file's virtual dataset; when there is none, the
  image files are found by listing the scan folder (top level only).
\end{itemize}
\end{hyp}

\subsection{Module header and helpers}

\pycode{bliss_hdf5_reader}{_header}

The constants list the positioner names that are time axes, not motors (\fn{TIME_AXES}),
the commands without a motor (\fn{TIME_SCANS}) and the single-point commands.

\api[\texttt{\_open}]{\texttt{\_open(path)}}
Opens a file read-only with HDF5 file locking disabled (\fn{locking=False}), so that a file
that BLISS is still writing can be read at the beamline; falls back to the default for
h5py older than 3.5, which has no \fn{locking} argument. Returns an \fn{h5py.File}, to be
used as a context manager.
\pycode{bliss_hdf5_reader}{_open}

\api[\texttt{\_to\_str}, \texttt{\_read\_scalar}, \texttt{\_is\_number}]{\texttt{\_to\_str(value)}, \texttt{\_read\_scalar(group, name)}, \texttt{\_is\_number(token)}}
\fn{_to_str} decodes bytes (UTF-8, with replacement of invalid sequences) and zero-dimensional
arrays to a Python string. \fn{_read_scalar} returns \fn{group[name]} as a Python scalar or
string, or \fn{None} if the dataset is missing. \fn{_is_number} tells whether a title token
parses as a float.
\pycode{bliss_hdf5_reader}{_to_str}
\pycode{bliss_hdf5_reader}{_read_scalar}
\pycode{bliss_hdf5_reader}{_is_number}

\api[\texttt{\_parse\_key}]{\texttt{\_parse\_key(scan, subscan)}}
Accepts a scan given as an integer (\fn{12}), a string (\fn{"12"}) or a full entry key
(\fn{"12.2"}). Returns \fn{(number, subscan, key)}; the subscan of a full key takes
precedence over the \fn{subscan} argument.
\pycode{bliss_hdf5_reader}{_parse_key}

\api[\texttt{list\_scans}]{\texttt{list\_scans(h5\_filepath)}}
Lists the scans of a master file without reading any data.
\blk{Returns} a list of dictionaries, one per top-level group, with keys \fn{key},
\fn{number} (integer part of the key, or \fn{None}), \fn{title}, \fn{npoints} (planned, from
\fn{scan_parameters}, if recorded), \fn{end_reason} and \fn{start_time}, sorted numerically by
scan then subscan (so \fn{"10.1"} comes after \fn{"9.1"}).
\pycode{bliss_hdf5_reader}{list_scans}

\api{\texttt{ScanMotor}}
A motor driven by the scan: \fn{name}, \fn{dim} (the scan dimension it drives, 0 = slowest),
\fn{npoints} along that dimension, \fn{start} and \fn{stop} as written in the title
(relative to the current position for d-scans, flagged by \fn{relative}). Its \fn{step}
property is $(\text{stop}-\text{start})/(n-1)$, or \fn{None} when undefined. Several motors can
share a dimension (\fn{a2scan}, the slow motors of an inferred mesh).
\pycode{bliss_hdf5_reader}{ScanMotor}
\pycode{bliss_hdf5_reader}{ScanMotor.step}

\api{\texttt{FrameSource}}
Location of one detector frame on disk: the source \fn{file}, the internal \fn{dataset}
path, and the frame \fn{index} in that dataset. Produced by \fn{BlissScan.frame_sources}.
\pycode{bliss_hdf5_reader}{FrameSource}

\subsection{Title parsers}

Each parser receives the command and the list of title tokens (commas are treated as
spaces) and returns \fn{(shape, motors, count_time)}, or \fn{None} if the tokens do not
follow the expected grammar---in which case the geometry is inferred from the data.

\api[\texttt{\_parse\_step\_scan}]{\texttt{\_parse\_step\_scan(cmd, args)}}
Step scans with $N$ motors: $3N+2$ tokens, each triple a motor name and two numbers, then an
integer number of intervals and a count time. The number of points is intervals $+1$.
Commands starting with \fn{d}, and \fn{lineup}, are relative.
\pycode{bliss_hdf5_reader}{_parse_step_scan}

\api[\texttt{\_parse\_mesh}]{\texttt{\_parse\_mesh(cmd, args)}}
\fn{amesh}/\fn{dmesh}: the first motor is the fast one, but the returned shape and motor
list put the slow motor first, so that every 2D scan has shape $(n_\mathrm{slow},
n_\mathrm{fast})$ whatever the command.
\pycode{bliss_hdf5_reader}{_parse_mesh}

\api[\texttt{\_parse\_fscan\_nd}]{\texttt{\_parse\_fscan\_nd(cmd, args)}}
Continuous scans \fn{fscan}, \fn{fscan2d}, \fn{fscan3d}: $d$ groups of (motor, start, step,
npoints), the \emph{last} motor being the fast one; the stop position is $\text{start} +
\text{step}\,(n-1)$.
\pycode{bliss_hdf5_reader}{_parse_fscan_nd}

\api[\texttt{\_parse\_fnscan}]{\texttt{\_parse\_fnscan(cmd, args)}}
Continuous scans of several motors along one dimension (\fn{f2scan}, \fn{f3scan}): the
number of motors is read from the command name.
\pycode{bliss_hdf5_reader}{_parse_fnscan}

\api[\texttt{\_parse\_time\_scan}, \texttt{\_parse\_single\_point}]{\texttt{\_parse\_time\_scan(cmd, args)}, \texttt{\_parse\_single\_point(cmd, args)}}
Scans without motors. A \fn{timescan} without a number of points runs until it is stopped:
its planned size is unknown (shape $(0,)$ at this stage, replaced by the recorded length).
\fn{ct}/\fn{sct} have shape $()$.
\pycode{bliss_hdf5_reader}{_parse_time_scan}
\pycode{bliss_hdf5_reader}{_parse_single_point}

\api[\texttt{\_parse\_title}]{\texttt{\_parse\_title(title)}}
Dispatches on the first token with regular expressions (\texttt{[ad]\textbackslash d*scan\textbar{}lineup},
\texttt{f\textbackslash d+scan}) and the command tables. Returns \fn{(command, parsed)} with \fn{parsed}
\fn{None} for unknown commands.
\pycode{bliss_hdf5_reader}{_parse_title}

\subsection{\texttt{BlissScan}}

\api[\texttt{BlissScan}]{\texttt{BlissScan(h5\_filepath, scan\_number, subscan=1, image\_prefix="eiger", image\_suffix=".h5", verbose=False)}}
One scan of a master file, of any type. Instantiation reads only metadata; data are read on
demand.
\begin{args}
\item[h5\_filepath] path of the master file.
\item[scan\_number] integer, or full key \fn{"12.1"}.
\item[subscan] used when \fn{scan_number} is an integer (default 1).
\item[image\_prefix] preferred image detector: the first image dataset whose name starts
  with it (default \fn{"eiger"}); \fn{None} takes the first image detector.
\item[image\_suffix] extension of image files when they have to be found by listing the
  scan folder.
\item[verbose] print the summary after \fn{read_data()}.
\end{args}
\blk{Main attributes} \fn{title}, \fn{command} (and its former name \fn{scan_type}),
\fn{shape}, \fn{motors} (list of \fn{ScanMotor}), \fn{npoints} (recorded),
\fn{npoints_planned}, \fn{count_time}, \fn{snake}, \fn{geometry_source} (\fn{"title"} or
\fn{"data"}), \fn{start_time}, \fn{end_time}, \fn{end_reason}, \fn{writer_status},
\fn{sample_name}, \fn{scan_folder}; after \fn{read_data()}: \fn{counters} (raw arrays),
\fn{grids} (scalar counters on the scan grid), \fn{image_path_list}, \fn{nb_images},
\fn{nb_frames}.
\pycode{bliss_hdf5_reader}{BlissScan}
\pycode{bliss_hdf5_reader}{BlissScan.__init__}

\api[\texttt{BlissScan.\_read\_metadata}]{\texttt{BlissScan.\_read\_metadata()}}
The core of the reader. It opens the file once and reads the entry's scalars, the shapes and
types of every measurement dataset (not the data), the positioner names
(\fn{_find_positioners}), the number of recorded points (\fn{_count_acquired_points}) and the
positions of every positioner recorded once per point (kept in a cache for geometry
inference and snake detection). It then:
\begin{enumerate}[nosep]
\item parses the title; if that fails, infers the geometry from the motor positions;
\item takes the count time from \fn{scan_parameters}, else from the title;
\item determines the planned number of points: the product of the title shape, overridden
  by \fn{scan_parameters/npoints} if they disagree (with a warning);
\item for 0D and 1D scans, replaces the shape by what was recorded;
\item detects back-and-forth meshes.
\end{enumerate}
Raises \fn{FileNotFoundError} or \fn{KeyError} if the file or the scan does not exist.
\pycode{bliss_hdf5_reader}{BlissScan._read_metadata}

\api[\texttt{BlissScan.\_find\_positioners}]{\texttt{BlissScan.\_find\_positioners(entry, meas)} (static)}
Motors among the measurement datasets: names present in \fn{positioners_start}, and datasets
whose link target is an \fn{instrument/} group of class \fn{NXpositioner}. The second rule
catches the fast motor of continuous scans, recorded as an encoder channel. Time axes are
excluded.
\pycode{bliss_hdf5_reader}{BlissScan._find_positioners}

\api[\texttt{BlissScan.\_count\_acquired\_points}]{\texttt{BlissScan.\_count\_acquired\_points(npts\_param)}}
The number of recorded points is the most common first-dimension length among the
measurement datasets (ties broken towards the larger value). Lengths above the planned number
of points are discarded first, which removes the $2n+1$ ``\fn{*_raw}'' channels of fast
scans.
\pycode{bliss_hdf5_reader}{BlissScan._count_acquired_points}

\api[\texttt{BlissScan.\_infer\_geometry}]{\texttt{BlissScan.\_infer\_geometry(npts\_planned=None)}}
Geometry for unrecognised titles. Motors that do not move (peak-to-peak range below
$10^{-9}$ of their magnitude) are ignored. If the moving motors form a mesh
(\fn{_infer_mesh}), the shape is $(n_\mathrm{slow}, n_\mathrm{fast})$, extended to the planned
number of rows for an aborted mesh; otherwise the scan is 1D and every moving motor drives
dimension 0.
\pycode{bliss_hdf5_reader}{BlissScan._infer_geometry}

\api[\texttt{BlissScan.\_row\_length}]{\texttt{BlissScan.\_row\_length(pos)} (static)}
Tests whether a position series is the slow motor of a mesh: it must jump at regular
intervals (one jump per row, all rows of equal length, at least two full rows) and stay
constant within a row (range within a row below a tenth of the median jump). Returns the row
length or \fn{None}.
\pycode{bliss_hdf5_reader}{BlissScan._row_length}

\api[\texttt{BlissScan.\_infer\_mesh}]{\texttt{BlissScan.\_infer\_mesh(moving)} (class method)}
For each candidate row length (shortest first), the slow motors are those with that row
length, and the fast motors those that repeat the same positions on the second row as on
the first, in the same or the reversed order (mismatch below a tenth of their range).
Returns \fn{((n_slow, n_fast), motors)} or \fn{None}.
\pycode{bliss_hdf5_reader}{BlissScan._infer_mesh}

\api[\texttt{BlissScan.\_detect\_snake}]{\texttt{BlissScan.\_detect\_snake()}}
A 2D scan is back-and-forth if the second row of the fast motor is closer to the reversed
first row than to the first row itself.
\pycode{bliss_hdf5_reader}{BlissScan._detect_snake}

\api[Properties: \texttt{ndim}, \texttt{is\_complete}, \texttt{motor\_names}, \texttt{positioner\_names}; \texttt{motors\_of\_dim}]{Properties: \texttt{ndim}, \texttt{is\_complete}, \texttt{motor\_names}, \texttt{positioner\_names}; \texttt{motors\_of\_dim(dim)}}
\fn{is_complete} is \fn{end_reason == "SUCCESS"} when the end reason is recorded, otherwise
a comparison of recorded and planned points. \fn{positioner_names} lists every motor
recorded at each point, scanned or not.
\pycode{bliss_hdf5_reader}{BlissScan.ndim}
\pycode{bliss_hdf5_reader}{BlissScan.is_complete}
\pycode{bliss_hdf5_reader}{BlissScan.motor_names}
\pycode{bliss_hdf5_reader}{BlissScan.positioner_names}
\pycode{bliss_hdf5_reader}{BlissScan.motors_of_dim}

\api[\texttt{BlissScan.index\_grid}]{\texttt{BlissScan.index\_grid} (property)}
The flat point index at each grid position, $-1$ where no point was recorded; odd rows are
reversed for a back-and-forth mesh. It is the single object that places per-point data on
the grid, used by \fn{to_grid} and by the image scripts.
\pycode{bliss_hdf5_reader}{BlissScan.index_grid}

\api[\texttt{get\_available\_counters}, \texttt{get}, \texttt{\_\_getitem\_\_}, \texttt{positions}]{\texttt{get\_available\_counters()} (alias \texttt{list\_counters}), \texttt{get(name, grid=False)}, \texttt{\_\_getitem\_\_(name)}, \texttt{positions(motor, grid=False)}}
\fn{get_available_counters} returns \texttt{\{name: \{"shape", "dtype"\}\}} for every
measurement dataset without reading data. \fn{get} reads one dataset (or returns it from
\fn{counters} if \fn{read_data} loaded it), warns when it would load more than 2\,GB, and
with \fn{grid=True} reshapes the point axis to the scan shape: an MCA dataset of a mesh
becomes $(n_\mathrm{slow}, n_\mathrm{fast}, n_\mathrm{ch})$. \fn{scan["name"]} is
\fn{get(name, grid=True)}; \fn{positions} returns floats.
\pycode{bliss_hdf5_reader}{BlissScan.get_available_counters}
\pycode{bliss_hdf5_reader}{BlissScan.get}
\pycode{bliss_hdf5_reader}{BlissScan.__getitem__}
\pycode{bliss_hdf5_reader}{BlissScan.positions}

\api[\texttt{BlissScan.to\_grid}]{\texttt{BlissScan.to\_grid(data, name="array")}}
Places per-point data on the scan grid. Complete, non-snake data are simply reshaped. With
missing points (aborted mesh) the output is filled with NaN (integer data are promoted to
floating point for that) and the recorded points are placed through \fn{index_grid}, which
also un-snakes back-and-forth rows. A counter with one value more than the grid (BLISS
sometimes records one point ahead on abort) is truncated. Raises \fn{ValueError} when the
first axis is not one value per point.
\pycode{bliss_hdf5_reader}{BlissScan.to_grid}

\api[\texttt{BlissScan.read\_data}]{\texttt{BlissScan.read\_data(extra\_counters=None)}}
Loads the scanned motors and the requested counters into \fn{counters}, grids every 1D
counter into \fn{grids}, and locates the image files (\fn{_process_images}). Missing names
produce warnings, not errors.
\pycode{bliss_hdf5_reader}{BlissScan.read_data}

\subsection{Detector frames}

\api[\texttt{image\_detectors}, \texttt{spectrum\_detectors} ; \texttt{\_pick\_detector}]{\texttt{image\_detectors}, \texttt{spectrum\_detectors} (properties); \texttt{\_pick\_detector(detector)}}
Datasets with one 2D image per point (3D datasets) and with one spectrum per point (2D
datasets). \fn{_pick_detector} returns the named detector, or the first image detector
matching \fn{image_prefix}, or the first image detector.
\pycode{bliss_hdf5_reader}{BlissScan.image_detectors}
\pycode{bliss_hdf5_reader}{BlissScan.spectrum_detectors}
\pycode{bliss_hdf5_reader}{BlissScan._pick_detector}

\api[\texttt{BlissScan.frame\_sources}]{\texttt{BlissScan.frame\_sources(detector=None)}}
Reads the mapping of the virtual dataset: for each virtual source, the selection bounds in
the virtual and source spaces give the range of frames it holds; relative file names are
resolved against the master file's folder. Returns one \fn{FrameSource} per frame, or an
empty list when the detector is not a virtual dataset.
\pycode{bliss_hdf5_reader}{BlissScan.frame_sources}

\api[\texttt{read\_frame}, \texttt{iter\_frames}]{\texttt{read\_frame(index, detector=None, roi=None)}, \texttt{iter\_frames(detector=None, indices=None, roi=None)}}
One frame by point index, read through the master file (HDF5 resolves the virtual mapping),
optionally restricted to \fn{roi = (row slice, column slice)} so that only that part is read
from disk. \fn{iter_frames} yields frames while keeping the file open, which is much faster
than repeated \fn{read_frame} calls.
\pycode{bliss_hdf5_reader}{BlissScan.read_frame}
\pycode{bliss_hdf5_reader}{BlissScan.iter_frames}

\api[\texttt{BlissScan.\_process\_images}]{\texttt{BlissScan.\_process\_images()}}
Fills \fn{image_path_list} with the distinct image files in acquisition order (from the
virtual dataset, or by listing \fn{scanNNNN/<prefix>*<suffix>}), \fn{nb_images} (files) and
\fn{nb_frames} (frames).
\pycode{bliss_hdf5_reader}{BlissScan._process_images}

\subsection{Other methods, the mesh subclass and the command line}

\api[\texttt{positioners\_start}, \texttt{get\_relative\_position\_grid}, \texttt{summary}, \texttt{get\_summary}, \texttt{\_\_repr\_\_}]{\texttt{positioners\_start()}, \texttt{get\_relative\_position\_grid(motor\_name, scale=1000.0)}, \texttt{summary()}, \texttt{get\_summary()}, \texttt{\_\_repr\_\_()}}
\fn{positioners_start} returns the snapshot of every motor at the start of the scan.
\fn{get_relative_position_grid} returns the gridded positions minus that of the first
recorded point, times \fn{scale} (mm to \textmu m by default). \fn{summary} is a
human-readable description: shape and its source, points recorded and planned, completion,
snake, count time, motors with their ranges and steps, detectors, loaded counters.
\pycode{bliss_hdf5_reader}{BlissScan.positioners_start}
\pycode{bliss_hdf5_reader}{BlissScan.get_relative_position_grid}
\pycode{bliss_hdf5_reader}{BlissScan.summary}
\pycode{bliss_hdf5_reader}{BlissScan.get_summary}
\pycode{bliss_hdf5_reader}{BlissScan.__repr__}

\api{\texttt{BlissMeshScan}}
A \fn{BlissScan} that must be two-dimensional (\fn{ValueError} otherwise), kept for scripts
written for the former mesh-only reader. It adds the attributes \fn{slow_/fast_motor_name},
\fn{_start}, \fn{_stop}, \fn{_steps}, \fn{_intervals}, \fn{_step_size}, \fn{scan_shape} and
\fn{expected_points}.
\pycode{bliss_hdf5_reader}{BlissMeshScan}
\pycode{bliss_hdf5_reader}{BlissMeshScan._read_metadata}

\api[\texttt{main}]{\texttt{main()}}
Command line: \fn{python bliss_hdf5_reader.py file.h5} lists the scans (key, planned
points, end reason, title); with scan keys after the file name it prints their summaries and
dataset shapes.
\pycode{bliss_hdf5_reader}{main}
